\documentclass[12pt]{article}
\usepackage[utf8]{inputenc}
\usepackage[T1]{fontenc}
\usepackage{hyperref}
\title{Ngen Language Resources for K-12 Curriculum Development}
\author{Compiled by Claude Sonnet 5.5 (claude-sonnet-5-5) for Andr\'as Kornai}
\date{}
\begin{document}
\maketitle

\section{Introduction}
Ngen (ISO 639-3: ssv) is an endangered Austronesian language of Vanuatu. In the search results it is placed among the East Santo languages on Espiritu Santo. Speaker counts differ between sources: Wikipedia gives about 450 and Verbix lists 800. Ngen should not be confused with Nguna (North Efate, ISO 639-3 llp). The search found very little material on Ngen, and none of it is educational.

\section{Primary Educational Resources}
No primers, textbooks or school programmes for Ngen turned up in the search. The language is not known to be taught in schools.

\section{Linguistic and Reference Materials}
\begin{itemize}
\item Wikipedia, ``Ngen language'': a short overview with classification and speaker estimates. \url{https://en.wikipedia.org/wiki/Ngen_language}
\item Verbix, endangered languages page for Ngen: a short descriptive entry. \url{https://docs.verbix.com/EndangeredLanguages/Ngen}
\end{itemize}
The search found no grammar or dictionary of Ngen.

\section{Religious and Community Resources}
None were found in the search.

\section{Digital and Audio Resources}
None were found in the search.

\section{Curriculum Development Context}
Almost nothing is documented for Ngen. A curriculum would therefore need new fieldwork and community-led documentation, and a basic orthography and word lists would have to come first. Any plan should be checked with Ngen speakers and with Vanuatu educational authorities.

\end{document}